\section{The \texttt{Source\_custom} McStas Component}
A flexible pulsed or continuous source with customisable parameters. Multiple pulses can be simulated over time.

\subsection*{Identification}
\begin{itemize}
  \item \textbf{Author:} Pablo Gila-Herranz, based on 'Source\_pulsed' by K. Lieutenant
  \item \textbf{Origin:} Materials Physics Center (CFM-MPC), CSIC-UPV/EHU
  \item \textbf{Date:} May 2025
\end{itemize}

\subsection*{Description}
Produces a custom pulse spectrum with a wavelength distribution as a sum of up to 3 Maxwellian distributions and one of undermoderated neutrons.

\textbf{Usage:}

By default, all Maxwellian distributions are assumed to come from anywhere in the moderator. A custom radius r\_i can be defined such that the central circle is at temperature T1, and the surrounding ring is at temperature T2. The third Maxwellian distribution at temperature T3 always comes from anywhere in the moderator.

Multiple pulses can be simulated over time with the n\_pulses parameter. The input flux is in units of [1/(cm\textasciicircum{}2 sr s \AA{})], and the output flux is per second, regardless of the number of pulses. This means that the intensity spreads over n\_pulses, so that the units of neutrons per second are preserved. To simulate only the neutron counts of a single pulse, the input flux should be divided by the frequency. Bear in mind that the maximum emission time is given by t\_pulse * tmax\_multiplier, so short pulses with long tails may require a higher tmax\_multiplier value (3 by default).

To simulate a continuous source, set a pulse length equal to the pulse period (e.g. t\_pulse=1.0, freq=1.0). Note that this continuous beam is simulated over time, unlike other continuous sources in McStas which produce a Dirac delta at time 0.

Parameters xwidth and yheight must be provided for rectangular moderators, otherwise a circular shape is assumed.

\textbf{Model description:}

The normalised Maxwellian distribution for moderated neutrons is defined by [1] \textless{}div class="latex"\textgreater{} \$M(\textbackslash{}lambda)=\textbackslash{}frac\{2a\textasciicircum{}2\}\{T\textasciicircum{}2\textbackslash{}lambda\textasciicircum{}5\}\textbackslash{}exp\textbackslash{}left(-\textbackslash{}frac\{a\}\{T\textbackslash{}lambda\textasciicircum{}\{2\}\}\textbackslash{}right)\$ \textless{}/div\textgreater{} where \textless{}div class="latex"\textgreater{} \$a=\textbackslash{}left(\textbackslash{}frac\{h\textasciicircum{}2\}\{2m\_\{N\}k\_\{B\}\}\textbackslash{}right)\$ \textless{}/div\textgreater{} and the joining function for the under-moderated neutrons is given by [1] \textless{}div class="latex"\textgreater{} \$M(\textbackslash{}lambda)\_\{um\}=\textbackslash{}frac\{1\}\{\textbackslash{}lambda(1+\textbackslash{}exp(\textbackslash{}lambda\textbackslash{}chi-\textbackslash{}kappa))\}\$ \textless{}/div\textgreater{}

The normalised time structure of the long pulse is defined by [2] \textless{}div class="latex"\textgreater{} \$N\_\{t\textless{}=t\_p\}=1-\textbackslash{}exp\textbackslash{}left(-\textbackslash{}frac\{t\}\{\textbackslash{}tau/n\}\textbackslash{}right)\$ \textless{}/div\textgreater{} \textless{}div class="latex"\textgreater{} \$N\_\{t\textgreater{}t\_p\}=\textbackslash{}exp\textbackslash{}left(-\textbackslash{}frac\{t-t\_p\}\{\textbackslash{}tau\}\textbackslash{}right)-\textbackslash{}exp\textbackslash{}left(-\textbackslash{}frac\{t\}\{\textbackslash{}tau/n\}\textbackslash{}right)\$ \textless{}/div\textgreater{} where tp is the pulse period, tau is the pulse decay time constant, and n is the ratio of decay to ascend time constants.

\textbf{Parameters for some sources:}

HBS thermal source: t\_pulse=0.016, freq=24.0, xwidth=0.04, yheight=0.04,

\begin{verbatim}
T1=325.0, I1=0.68e+12, tau1=0.000125,
I_um=2.47e+10, chi_um=2.5
\end{verbatim}

HBS cold source: t\_pulse=0.016, freq=24.0, radius=0.010,

\begin{verbatim}
T1=60.0, I1=1.75e+12, tau1=0.000170,
I_um=3.82e+10, chi_um=0.9
\end{verbatim}

HBS bi-spectral:

\begin{verbatim}
t_pulse=0.016, freq=24.0, radius=0.022, r_i=0.010,
T1= 60.0, I1=1.75e+12, tau1=0.000170,
T2=305.0, I2=0.56e+12, tau2=0.000130,
I_um=3.82e+10, chi_um=2.5
\end{verbatim}

PSI cold source: t\_pulse=1, freq=1, xwidth=0.1, yheight=0.1,

\begin{verbatim}
T1=296.2, I1=8.5e+11, T2=40.68, I2=5.2e+11
\end{verbatim}

\textbf{References:} [1] J. M. Carpenter et al, Nuclear Instruments and Methods in Physics Research Section A, 234, 542-551 (1985). [2] J. M. Carpenter and W. B. Yelon, Methods in Experimental Physics 23, p. 127, Chapter 2, Neutron Sources (1986).

\subsection*{Input parameters}
Parameters in \textbf{boldface} are required; the others are optional.

\begin{longtable}{p{0.22\textwidth}p{0.12\textwidth}p{0.46\textwidth}p{0.14\textwidth}}
\toprule
\textbf{Name} & \textbf{Unit} & \textbf{Description} & \textbf{Default} \\
\midrule
\endhead
target\_index & 1 & Relative index of component to focus at, e.g. next is +1 (used to compute 'dist' automatically) & 1 \\
dist & m & Distance from the source to the target & 0.0 \\
focus\_xw & m & Width of the target (focusing rectangle) & 0.02 \\
focus\_yh & m & Height of the target (focusing rectangle) & 0.02 \\
xwidth & m & Width of the source & 0.0 \\
yheight & m & Height of the source & 0.0 \\
radius & m & Outer radius of the source (ignored if xwidth and yheight are set) & 0.010 \\
r\_i & m & Radius of a central circle at temperature T1, sorrounded by a ring at temperature T2 (ignored by default) & 0.0 \\
\textbf{Lmin} & \AA{} & Lower edge of the wavelength distribution &  \\
\textbf{Lmax} & \AA{} & Upper edge of the wavelength distribution &  \\
\textbf{freq} & Hz & Frequency of pulses &  \\
\textbf{t\_pulse} & s & Proton pulse length &  \\
tmax\_multiplier & 1 & Defined maximum emission time at moderator (tmax = tmax\_multiplier * t\_pulse); 1 for continuous sources & 3.0 \\
n\_pulses & 1 & Number of pulses to be simulated (ignored for continuous sources) & 1 \\
T1 & K & Temperature of the 1st Maxwellian distribution; for r\_i \textgreater{} 0 it is only used for the central radii r \textless{} r\_i & 0.0 \\
I1 & 1/(cm$^{2}$ sr s \AA{}) & Flux per solid angle of the 1st Maxwellian distribution (integrated over the whole wavelength range) & 0.0 \\
tau1 & s & Pulse decay time constant of the 1st Maxwellian distribution & 0.000125 \\
T2 & K & Temperature of the 2nd Maxwellian distribution (0 to disable); for r\_i \textgreater{} 0 it is only used for the external ring r \textgreater{} r\_i & 0.0 \\
I2 & 1/(cm$^{2}$ sr s \AA{}) & Flux per solid angle of the 2nd Maxwellian distribution & 0.0 \\
tau2 & s & Pulse decay time constant of the 2nd Maxwellian distribution & 0.000125 \\
T3 & K & Temperature of the 3rd Maxwellian distribution (0 to disable) & 0.0 \\
I3 & 1/(cm$^{2}$ sr s \AA{}) & Flux per solid angle of the 3rd Maxwellian distribution & 0.0 \\
tau3 & s & Pulse decay time constant of the 3rd Maxwellian distribution & 0.000125 \\
n\_mod & 1 & Ratio of pulse decay constant to pulse ascend constant of moderated neutrons (n = tau\_decay / tau\_ascend) & 1 \\
I\_um & 1/(cm$^{2}$ sr s \AA{}) & Flux per solid angle for the under-moderated neutrons & 0.0 \\
tau\_um & s & Pulse decay time constant of under-moderated neutrons & 0.000012 \\
n\_um & 1 & Ratio of pulse decay constant to pulse ascend constant of under-moderated neutrons & 1 \\
chi\_um & 1/\AA{} & Factor for the wavelength dependence of under-moderated neutrons & 2.5 \\
kap\_um & 1 & Scaling factor for the flux of under-moderated neutrons & 2.2 \\
\bottomrule
\end{longtable}

\subsection*{Links}
\begin{itemize}
  \item Component source code found in file \texttt{Source\_custom.comp}.
  \item This model is implemented in an \htmladdnormallink{interactive notebook}{https://github.com/pablogila/Source\_custom} to fit custom neutron sources.
\end{itemize}
\IfFileExists{contrib/Source_custom_static.tex}{\input{contrib/Source_custom_static.tex}}{}